\section{Introduction}

Multi-Agent Debate (MAD) improves the accuracy of Large Language Model (LLM) agents across diverse domains, including research, mathematics, and coding~\citep{Du2024Improving, Li2023Camel, Wu2024Autogen}. By exchanging intermediate reasoning steps and answers across rounds, agents can revise their responses based on peer information and achieve accuracy gains beyond single-agent settings. However, MAD becomes costly as the communication graph grows. Each shared message is appended to another agent’s context, increasing the input-token cost and KV cache usage. With more agents or debate rounds, this overhead compounds quickly, limiting scalability.

Recent work reduces the cost of MAD by removing nodes (i.e., participating agents) and/or edges (i.e., communication links between agents) from all-to-all communication graphs~\citep{Liu2025Groupdebate, Chen2025Sid, Zeng2025S2Mad, Zhu2026Demystifying}. These methods typically leverage information available prior to the debate (e.g., token-level log-likelihoods, LLM's self-stated confidence) to identify and remove potentially low-utility communications. For example, outputs with high min. log-likelihood values are often considered to be reliable, and therefore are likely to be correct without further debate~\citep{Chen2025Sid}.

In this work, we uncover that although pre-debate signals are effective in common cases, their informativeness can collapse on challenging problem instances relative to the LLM’s capabilities. On such instances, incorrect agents may hallucinate with high confidence, while correct agents may remain uncertain~\citep{Kakai2025Why}. Methods based on prior signals can prematurely exclude agents or communication links that would otherwise be beneficial, thereby reducing MAD's accuracy gains.

To address this limitation, we take a Bayesian-inspired view of MAD: prior agent signals can be unreliable under hallucination, whereas debate outcomes provide posterior evidence of reliability. We focus on whether an agent preserves its answer after incorporating peer messages. Intuitively, an agent supported by sound reasoning is less likely to be convinced by flawed peer arguments. We propose \textit{Survival Rate} (SVR), which measures how often an agent retains its original belief after being challenged. Because SVR is grounded in observed post-debate behavior, it provides a more hallucination-robust signal for identifying trustworthy reasoning processes and answers.

We propose \sysname, an efficient MAD framework that incrementally builds a communication graph using posterior debate outcomes. \sysname initializes agent correctness scores from prior signals, probes $S$ communication links to high-prior receivers, and updates scores based on whether agents retain or revise their answers. It then prioritizes high-score agents in later communications and terminates once an agent’s score exceeds a threshold. This greedy strategy identifies reliable solutions early, reducing token costs. Across two LLMs and datasets, \sysname reduces token costs by up to 61\% while matching or improving accuracy relative to the most accurate competing MAD baseline.
